\documentclass[10pt,a4paper]{amsart}
\usepackage{amsmath,amssymb,amsthm,graphicx}
% Retained preamble for the Stacks Project chapter {CH} (shared project preamble).
% The chapter file pinned-original-source.tex only carries \input{{preamble}} on its first line,
% so no preamble text lives inside the chapter itself.  The Stacks Project uses one single shared
% preamble for every chapter: the snapshot stages it three times (sources/stacks-chow/preamble.tex,
% sources/stacks-etale-cohomology/preamble.tex and focused-seeds/AGP/sources/stacks-snapshot-20261011/
% preamble.tex) and those copies are identical in size.  The block below is copied mechanically,
% line for line, from sources/stacks-chow/preamble.tex lines 14-18 and 154-217: the Xy-pic setup and
% the native theorem declarations together with the mathematical macros.  Nothing was added, removed
% or redefined; the wrapper class and the amsmath/amssymb/amsthm/graphicx base packages are supplied
% by the assembler.
% For commutative diagrams we use Xy-pic
% Chapter: stacks-spaces-more-morphisms.tex -- More on Morphisms of Spaces.
\usepackage[all]{xy}

% We use 2cell for 2-commutative diagrams.
\xyoption{2cell}

% Theorem environments.
%
\theoremstyle{plain}
\newtheorem{theorem}[subsection]{Theorem}
\newtheorem{proposition}[subsection]{Proposition}
\newtheorem{lemma}[subsection]{Lemma}

\theoremstyle{definition}
\newtheorem{definition}[subsection]{Definition}
\newtheorem{example}[subsection]{Example}
\newtheorem{exercise}[subsection]{Exercise}
\newtheorem{situation}[subsection]{Situation}

\theoremstyle{remark}
\newtheorem{remark}[subsection]{Remark}
\newtheorem{remarks}[subsection]{Remarks}

\numberwithin{equation}{subsection}

% Macros
%
\def\lim{\mathop{\mathrm{lim}}\nolimits}
\def\colim{\mathop{\mathrm{colim}}\nolimits}
\def\Spec{\mathop{\mathrm{Spec}}}
\def\Hom{\mathop{\mathrm{Hom}}\nolimits}
\def\Ext{\mathop{\mathrm{Ext}}\nolimits}
\def\SheafHom{\mathop{\mathcal{H}\!\mathit{om}}\nolimits}
\def\SheafExt{\mathop{\mathcal{E}\!\mathit{xt}}\nolimits}
\def\Sch{\mathit{Sch}}
\def\Mor{\mathop{\mathrm{Mor}}\nolimits}
\def\Ob{\mathop{\mathrm{Ob}}\nolimits}
\def\Sh{\mathop{\mathit{Sh}}\nolimits}
\def\NL{\mathop{N\!L}\nolimits}
\def\CH{\mathop{\mathrm{CH}}\nolimits}
\def\proetale{{pro\text{-}\acute{e}tale}}
\def\etale{{\acute{e}tale}}
\def\QCoh{\mathit{QCoh}}
\def\Ker{\mathop{\mathrm{Ker}}}
\def\Im{\mathop{\mathrm{Im}}}
\def\Coker{\mathop{\mathrm{Coker}}}
\def\Coim{\mathop{\mathrm{Coim}}}

% Boxtimes
%
\DeclareMathSymbol{\boxtimes}{\mathbin}{AMSa}{"02}

%
% Macros for moduli stacks/spaces
%
\def\QCohstack{\mathcal{QC}\!\mathit{oh}}
\def\Cohstack{\mathcal{C}\!\mathit{oh}}
\def\Spacesstack{\mathcal{S}\!\mathit{paces}}
\def\Quotfunctor{\mathrm{Quot}}
\def\Hilbfunctor{\mathrm{Hilb}}
\def\Curvesstack{\mathcal{C}\!\mathit{urves}}
\def\Polarizedstack{\mathcal{P}\!\mathit{olarized}}
\def\Complexesstack{\mathcal{C}\!\mathit{omplexes}}
% \Pic is the operator that assigns to X its picard group, usage \Pic(X)
% \Picardstack_{X/B} denotes the Picard stack of X over B
% \Picardfunctor_{X/B} denotes the Picard functor of X over B
\def\Pic{\mathop{\mathrm{Pic}}\nolimits}
\def\Picardstack{\mathcal{P}\!\mathit{ic}}
\def\Picardfunctor{\mathrm{Pic}}
\def\Deformationcategory{\mathcal{D}\!\mathit{ef}}
\usepackage{cleveref}
\crefname{originalreference}{original reference}{original references}
\Crefname{originalreference}{Original reference}{Original references}
\makeatletter
\newlabel{decent-spaces-lemma-decent-finite-fibre}{{0AC7}{source}{}{sourcepointer-1}{}}
\newlabel{decent-spaces-lemma-decent-finite-fibre@cref}{{[originalreference][0][]0AC7}{[1][0][]source}{}{}{}}
\makeatother
\crefname{originalreference}{original reference}{original references}
\Crefname{originalreference}{Original reference}{Original references}
\makeatletter
\newlabel{decent-spaces-lemma-qf-and-qc-finite-fibre}{{0AC6}{source}{}{sourcepointer-2}{}}
\newlabel{decent-spaces-lemma-qf-and-qc-finite-fibre@cref}{{[originalreference][0][]0AC6}{[1][0][]source}{}{}{}}
\makeatother
\crefname{originalreference}{original reference}{original references}
\Crefname{originalreference}{Original reference}{Original references}
\makeatletter
\newlabel{decent-spaces-section-reasonable-decent}{{03I7}{source}{}{sourcepointer-3}{}}
\newlabel{decent-spaces-section-reasonable-decent@cref}{{[originalreference][0][]03I7}{[1][0][]source}{}{}{}}
\makeatother
\crefname{originalreference}{original reference}{original references}
\Crefname{originalreference}{Original reference}{Original references}
\makeatletter
\newlabel{groupoids-proposition-finite-flat-equivalence}{{03BM}{source}{}{sourcepointer-4}{}}
\newlabel{groupoids-proposition-finite-flat-equivalence@cref}{{[originalreference][0][]03BM}{[1][0][]source}{}{}{}}
\makeatother
\crefname{originalreference}{original reference}{original references}
\Crefname{originalreference}{Original reference}{Original references}
\makeatletter
\newlabel{more-morphisms-lemma-etale-splits-off-quasi-finite-part-technical}{{02LN}{source}{}{sourcepointer-5}{}}
\newlabel{more-morphisms-lemma-etale-splits-off-quasi-finite-part-technical@cref}{{[originalreference][0][]02LN}{[1][0][]source}{}{}{}}
\makeatother
\crefname{originalreference}{original reference}{original references}
\Crefname{originalreference}{Original reference}{Original references}
\makeatletter
\newlabel{morphisms-lemma-finite-proper}{{01WN}{source}{}{sourcepointer-6}{}}
\newlabel{morphisms-lemma-finite-proper@cref}{{[originalreference][0][]01WN}{[1][0][]source}{}{}{}}
\makeatother
\crefname{originalreference}{original reference}{original references}
\Crefname{originalreference}{Original reference}{Original references}
\makeatletter
\newlabel{morphisms-lemma-image-proper-is-proper}{{03GN}{source}{}{sourcepointer-7}{}}
\newlabel{morphisms-lemma-image-proper-is-proper@cref}{{[originalreference][0][]03GN}{[1][0][]source}{}{}{}}
\makeatother
\crefname{originalreference}{original reference}{original references}
\Crefname{originalreference}{Original reference}{Original references}
\makeatletter
\newlabel{spaces-lemma-space-presentation}{{0262}{source}{}{sourcepointer-8}{}}
\newlabel{spaces-lemma-space-presentation@cref}{{[originalreference][0][]0262}{[1][0][]source}{}{}{}}
\makeatother
\crefname{originalreference}{original reference}{original references}
\Crefname{originalreference}{Original reference}{Original references}
\makeatletter
\newlabel{spaces-morphisms-lemma-universally-closed-permanence}{{04NX}{source}{}{sourcepointer-9}{}}
\newlabel{spaces-morphisms-lemma-universally-closed-permanence@cref}{{[originalreference][0][]04NX}{[1][0][]source}{}{}{}}
\makeatother
\crefname{originalreference}{original reference}{original references}
\Crefname{originalreference}{Original reference}{Original references}
\makeatletter
\newlabel{spaces-properties-lemma-finite-fibres-presentation}{{03IJ}{source}{}{sourcepointer-10}{}}
\newlabel{spaces-properties-lemma-finite-fibres-presentation@cref}{{[originalreference][0][]03IJ}{[1][0][]source}{}{}{}}
\makeatother
\title{Splitting off the quasi-finite part after an etale neighbourhood, for algebraic spaces}
\author{The Stacks Project Authors}
\date{stacks-project:spaces-more-morphisms.tex:local-snapshot-20261011; GNU FDL 1.2 or later}
\begin{document}
\maketitle
\noindent\textit{Editorial note.} Counted span is the author's own proof environment, source lines 7494-7541, immediately following the statement at 7471-7492.\par\medskip
\noindent\textit{Editorial note.} Ten cross-chapter pointers (Decent Spaces, Groupoids, More on Morphisms, Morphisms, Spaces, Morphisms of Spaces) are recorded in assembly.external\_references.\par\medskip
\noindent\textit{Editorial note.} The proof's final steps (affineness of V from a finite flat equivalence relation, then finiteness of V -> U) are displayed inside the span.\par\medskip
\noindent\textit{Allowed standard prerequisites (editorial).}\begin{itemize}
\item The technical schemes-level splitting statement. More on Morphisms, lemma more-morphisms-lemma-etale-splits-off-quasi-finite-part-technical, cited in the proof. Supplies the scheme-level construction that is transported to algebraic spaces here.
\item Presentation of an algebraic space as a quotient by an etale equivalence relation. Spaces, lemma spaces-lemma-space-presentation, cited at proof lines 7528-7530. Realises the intermediate algebraic space V as a quotient of Z\_1 by R = Z\_1 x\_V Z\_1.
\item Affineness of a quotient by a finite flat equivalence relation. Groupoids, proposition groupoids-proposition-finite-flat-equivalence, cited at proof lines 7534-7535. Concludes that V is an affine scheme from finiteness of the projections of R.
\item A proper morphism with finite image is finite. Morphisms, lemmas morphisms-lemma-image-proper-is-proper and morphisms-lemma-finite-proper, cited at proof lines 7536-7539. Turns properness of V -> U into finiteness, which is the conclusion.
\item Finiteness of fibres and properness results for decent spaces. Decent Spaces lemmas (decent-finite-fibre, qf-and-qc-finite-fibre, section-reasonable-decent) and Morphisms of Spaces results cited in the proof. Control the finite-fibre hypotheses and properness in the algebraic-space setting.
\end{itemize}\medskip
\ifdefined\hypertarget\hypertarget{sourcepointer-1}{}\fi
\noindent\textit{Source reference.} Original reference 0AC7: focused-seeds/geometry/S9-spaces-etale-splits-off-quasi-finite-spaces/ref-tags.txt, source lines 5--5. Navigation pointer only: the counted proof cites this statement by \textbackslash{}ref as a standing input; the deduction the counted proof has to make is carried out inside the counted span, and no part of that deduction is delegated to this pointer.\par\medskip
\ifdefined\hypertarget\hypertarget{sourcepointer-2}{}\fi
\noindent\textit{Source reference.} Original reference 0AC6: focused-seeds/geometry/S9-spaces-etale-splits-off-quasi-finite-spaces/ref-tags.txt, source lines 6--6. Navigation pointer only: the counted proof cites this statement by \textbackslash{}ref as a standing input; the deduction the counted proof has to make is carried out inside the counted span, and no part of that deduction is delegated to this pointer.\par\medskip
\ifdefined\hypertarget\hypertarget{sourcepointer-3}{}\fi
\noindent\textit{Source reference.} Original reference 03I7: focused-seeds/geometry/S9-spaces-etale-splits-off-quasi-finite-spaces/ref-tags.txt, source lines 7--7. Navigation pointer only: the counted proof cites this statement by \textbackslash{}ref as a standing input; the deduction the counted proof has to make is carried out inside the counted span, and no part of that deduction is delegated to this pointer.\par\medskip
\ifdefined\hypertarget\hypertarget{sourcepointer-4}{}\fi
\noindent\textit{Source reference.} Original reference 03BM: focused-seeds/geometry/S9-spaces-etale-splits-off-quasi-finite-spaces/ref-tags.txt, source lines 8--8. Navigation pointer only: the counted proof cites this statement by \textbackslash{}ref as a standing input; the deduction the counted proof has to make is carried out inside the counted span, and no part of that deduction is delegated to this pointer.\par\medskip
\ifdefined\hypertarget\hypertarget{sourcepointer-5}{}\fi
\noindent\textit{Source reference.} Original reference 02LN: focused-seeds/geometry/S9-spaces-etale-splits-off-quasi-finite-spaces/ref-tags.txt, source lines 9--9. Navigation pointer only: the counted proof cites this statement by \textbackslash{}ref as a standing input; the deduction the counted proof has to make is carried out inside the counted span, and no part of that deduction is delegated to this pointer.\par\medskip
\ifdefined\hypertarget\hypertarget{sourcepointer-6}{}\fi
\noindent\textit{Source reference.} Original reference 01WN: focused-seeds/geometry/S9-spaces-etale-splits-off-quasi-finite-spaces/ref-tags.txt, source lines 10--10. Navigation pointer only: the counted proof cites this statement by \textbackslash{}ref as a standing input; the deduction the counted proof has to make is carried out inside the counted span, and no part of that deduction is delegated to this pointer.\par\medskip
\ifdefined\hypertarget\hypertarget{sourcepointer-7}{}\fi
\noindent\textit{Source reference.} Original reference 03GN: focused-seeds/geometry/S9-spaces-etale-splits-off-quasi-finite-spaces/ref-tags.txt, source lines 11--11. Navigation pointer only: the counted proof cites this statement by \textbackslash{}ref as a standing input; the deduction the counted proof has to make is carried out inside the counted span, and no part of that deduction is delegated to this pointer.\par\medskip
\ifdefined\hypertarget\hypertarget{sourcepointer-8}{}\fi
\noindent\textit{Source reference.} Original reference 0262: focused-seeds/geometry/S9-spaces-etale-splits-off-quasi-finite-spaces/ref-tags.txt, source lines 12--12. Navigation pointer only: the counted proof cites this statement by \textbackslash{}ref as a standing input; the deduction the counted proof has to make is carried out inside the counted span, and no part of that deduction is delegated to this pointer.\par\medskip
\ifdefined\hypertarget\hypertarget{sourcepointer-9}{}\fi
\noindent\textit{Source reference.} Original reference 04NX: focused-seeds/geometry/S9-spaces-etale-splits-off-quasi-finite-spaces/ref-tags.txt, source lines 13--13. Navigation pointer only: the counted proof cites this statement by \textbackslash{}ref as a standing input; the deduction the counted proof has to make is carried out inside the counted span, and no part of that deduction is delegated to this pointer.\par\medskip
\ifdefined\hypertarget\hypertarget{sourcepointer-10}{}\fi
\noindent\textit{Source reference.} Original reference 03IJ: focused-seeds/geometry/S9-spaces-etale-splits-off-quasi-finite-spaces/ref-tags.txt, source lines 14--14. Navigation pointer only: the counted proof cites this statement by \textbackslash{}ref as a standing input; the deduction the counted proof has to make is carried out inside the counted span, and no part of that deduction is delegated to this pointer.\par\medskip
\begin{lemma}
\label{lemma-etale-splits-off-quasi-finite-part}
Let $S$ be a scheme.
Let $f : X \to Y$ be a morphism of algebraic spaces over $S$.
Let $y \in |Y|$. Let $x_1, \ldots, x_n \in |X|$ mapping to $y$.
Assume that
\begin{enumerate}
\item $f$ is locally of finite type,
\item $f$ is separated,
\item $f$ is quasi-finite at $x_1, \ldots, x_n$, and
\item $f$ is quasi-compact or $Y$ is decent.
\end{enumerate}
Then there exists an \'etale morphism $(U, u) \to (Y, y)$
of pointed algebraic spaces and a decomposition
$$
U \times_Y X = W \amalg V
$$
into open and closed subspaces such that the morphism $V \to U$ is finite,
every point of the fibre of $|V| \to |U|$ over $u$ maps to an $x_i$,
and the fibre of $|W| \to |U|$ over $u$ contains no point mapping to an
$x_i$.
\end{lemma}

\begin{proof}
Let $(U, u) \to (Y, y)$ be an \'etale morphism of algebraic spaces
and consider the set of  $w \in |U \times_Y X|$ mapping to $u \in |U|$
and one of the $x_i \in |X|$. By
Decent Spaces, Lemma \ref{decent-spaces-lemma-qf-and-qc-finite-fibre}
(if $f$ is of finite type) or
Decent Spaces, Lemma \ref{decent-spaces-lemma-decent-finite-fibre}
(if $Y$ is decent) this set is finite.
It follows that we may replace $f$ by the base change
$U \times_Y X \to U$ and $x_1, \ldots, x_n$ by the set of these $w$.
In particular we may and do assume that $Y$ is an affine scheme,
whence $X$ is a separated algebraic space.

\medskip\noindent
Choose an affine scheme $Z$ and an \'etale morphism $Z \to X$ such that
$x_1, \ldots, x_n$ are in the image of $|Z| \to |X|$. The fibres of
$|Z| \to |X|$ are finite, see Properties of Spaces, Lemma
\ref{spaces-properties-lemma-finite-fibres-presentation}
(or the more general discussion in Decent Spaces, Section
\ref{decent-spaces-section-reasonable-decent}).
Let $\{z_1, \ldots, z_m\} \subset |Z|$ be the preimage of
$\{x_1, \ldots, x_n\}$. By More on Morphisms, Lemma
\ref{more-morphisms-lemma-etale-splits-off-quasi-finite-part-technical}
there exists an \'etale morphism $(U, u) \to (Y, y)$ such
that $U \times_Y Z = Z_1 \amalg Z_2$ with $Z_1 \to U$ finite and
$(Z_1)_y = \{z_1, \ldots, z_m\}$. We may assume that $U$ is affine
and hence $Z_1$ is affine too.

\medskip\noindent
Since $f$ is separated, the image $V$ of $Z_1 \to X$ is both open and closed
(Morphisms of Spaces, Lemma
\ref{spaces-morphisms-lemma-universally-closed-permanence}).
Set $W = X \setminus V$ to get a decomposition as in the lemma.
To finish the proof we have to show that $V \to U$ is finite.
As $Z_1 \to V$ is surjective and \'etale, $V$ is the quotient of
$Z_1$ by the \'etale equivalence relation $R = Z_1 \times_V Z_1$, see
Spaces, Lemma \ref{spaces-lemma-space-presentation}.
Since $f$ is separated, $V \to U$ is separated and $R$ is closed in
$Z_1 \times_U Z_1$. Since $Z_1 \to U$ is finite,
the projections $s, t : R \to Z_1$ are finite.
Thus $V$ is an affine scheme by
Groupoids, Proposition \ref{groupoids-proposition-finite-flat-equivalence}.
By Morphisms, Lemma \ref{morphisms-lemma-image-proper-is-proper}
we conclude that $V \to U$ is proper and by
Morphisms, Lemma \ref{morphisms-lemma-finite-proper}
we conclude that $V \to U$ is finite,
thereby finishing the proof.
\end{proof}

\end{document}
